%%%%%%%%%%%%%%%%%%%%%%%%%%%%%%%%%%%%%%%%%%%%%%%%%%%%%%%%%%%%%%%%%%
%% Mathematical Epidemiology: Survival Analysis of Erdos Problems
%%%%%%%%%%%%%%%%%%%%%%%%%%%%%%%%%%%%%%%%%%%%%%%%%%%%%%%%%%%%%%%%%%
\documentclass[11pt]{article}

\usepackage[margin=1in]{geometry}
\usepackage{amsmath,amssymb}
\usepackage{mathtools}
\usepackage{hyperref}
\usepackage{booktabs}
\usepackage{microtype}
\usepackage[numbers,sort&compress]{natbib}
\usepackage{xcolor}
\usepackage{enumitem}
\setlist{nosep}

\hypersetup{
  colorlinks=true,
  linkcolor=blue!60!black,
  citecolor=blue!60!black,
  urlcolor=blue!60!black
}

% Notation
\newcommand{\HR}{\mathrm{HR}}
\newcommand{\AF}{\mathrm{AF}}
\newcommand{\CI}{\mathrm{CI}}
\newcommand{\AIC}{\mathrm{AIC}}
\newcommand{\CIF}{\mathrm{CIF}}

\begin{document}

\title{Mathematical Epidemiology:\\
Survival Analysis of 1,183 Erd\H{o}s Problems}

\author{
  Alex Towell\thanks{Department of Computer Science,
  Southern Illinois University Edwardsville.
  Email: \texttt{atowell@siue.edu}.}
}

\date{}
\maketitle

\begin{abstract}
We present, to our knowledge, the first application of survival analysis
to the resolution dynamics of a mathematical problem corpus.
Treating 1,183 Erd\H{o}s problems as subjects in a right-censored
time-to-event study---with problem index serving as an ordinal time
proxy and resolution as the event of interest---we apply
Kaplan--Meier estimation, Cox proportional hazards regression,
competing risks analysis, and Weibull accelerated failure time
modeling.
Our principal findings are:
(i)~prize problems resolve significantly faster
($\HR = 2.60$, 95\%~$\CI$: $1.85$--$3.67$, $p < 0.001$);
(ii)~formalized problems exhibit 45\% lower hazard of resolution
($\HR = 0.55$, 95\%~$\CI$: $0.43$--$0.71$, $p < 0.001$),
consistent with a selection effect rather than a causal mechanism;
(iii)~the overall proved-to-disproved ratio is $2.63{:}1$,
but varies from $8.5{:}1$ in combinatorics to $1.3{:}1$ in set theory;
and (iv)~a Weibull AFT model with shape parameter $\hat\rho = 1.61$
indicates decreasing marginal hazard, with prize attachment
accelerating resolution by a factor of $0.60$ ($p < 0.001$).
The Cox model achieves a concordance index of $C = 0.60$ with a
significant global test ($\chi^2 = 65.4$, $\text{df} = 7$,
$p < 10^{-10}$).
\end{abstract}


%% ================================================================
\section{Introduction}
%% ================================================================

The Erd\H{o}s problem corpus---1,183 conjectures, questions, and
challenges posed by Paul Erd\H{o}s (1913--1996) and catalogued by
Tao et al.~\cite{erdosproblems}---provides an opportunity
to study the \emph{dynamics} of mathematical progress through a
quantitative lens. Each problem can be viewed as a patient entering
a longitudinal study at the time of its formulation, with resolution
(proof, disproof, or other definitive answer) constituting the event
of interest. Problems that remain open are right-censored.

While bibliometric approaches to mathematical productivity are
established~\cite{borjas2015,mihaljevi2019}, survival-analytic
methods have not previously been applied to a structured database of
open problems. The Erd\H{o}s corpus is particularly suited: problems
span six decades, carry metadata (tags, prizes, OEIS references,
formalization status), and exhibit a natural ordering by problem
number that serves as an ordinal time proxy.
We ask: which covariates predict faster resolution, how does the
hazard change across the corpus, and how do the competing risks
of proof versus disproof vary by subfield?

\smallskip\noindent
\textbf{Notation.}
Problem numbers $i \in \{1, \ldots, 1183\}$ serve as the
time variable~$T$. The event indicator is $\delta_i = 1$ if problem
$i$ has been resolved and $\delta_i = 0$ otherwise (right-censored).


%% ================================================================
\section{Data and Methods}
%% ================================================================

\subsection{Dataset}

The Erd\H{o}s Problems database~\cite{erdosproblems} records
$n = 1{,}183$ problems with the following fields: problem number,
status (open, proved, disproved, solved, falsifiable, verifiable,
decidable, independent, or Lean-verified variants), tags
(mathematical subfields), prize amount (USD or GBP), OEIS sequence
references, and formalization status. We restrict the primary analysis
to problems with a definitive binary outcome:
\[
  n_{\text{events}} = 483 \text{ (resolved)}, \quad
  n_{\text{censored}} = 650 \text{ (open)}, \quad
  n_{\text{total}} = 1{,}133.
\]
The remaining 50 observations with intermediate statuses
(falsifiable, verifiable, decidable, independent) are excluded
from the main survival models but included in the competing risks
analysis as a third event type.

Covariates are summarized in Table~\ref{tab:covariates}.

\begin{table}[ht]
\centering
\caption{Covariate definitions and descriptive statistics ($n = 1{,}133$).}
\label{tab:covariates}
\begin{tabular}{@{}llcc@{}}
\toprule
Covariate & Type & Prevalence/Mean & Events (\%) \\
\midrule
Prize ($> \$0$)       & Binary & $8.8\%$  ($100/1133$)   & $38.0\%$ \\
Formalized            & Binary & $33.4\%$ ($378/1133$)   & $21.7\%$ \\
OEIS sequences        & Binary & $25.2\%$ ($286/1133$)   & $26.2\%$ \\
Number of tags        & Count  & $\bar x = 1.63$         & --- \\
Tag solve rate        & Continuous $[0,1]$ & $\bar x = 0.40$ & --- \\
Number theory         & Binary & $47.7\%$ & $37.2\%$ \\
Graph theory          & Binary & $22.3\%$ & $49.0\%$ \\
Ramsey theory         & Binary & $9.3\%$  & $40.9\%$ \\
\bottomrule
\end{tabular}
\end{table}

\subsection{Statistical Methods}

\paragraph{Kaplan--Meier estimation.}
We estimate the survival function $\hat S(t)$ using the product-limit
estimator~\cite{kaplan1958}, with pointwise 95\% confidence intervals
via Greenwood's formula. Stratified curves are compared using the
log-rank test.

\paragraph{Cox proportional hazards regression.}
The hazard for problem $i$ is modeled as
\[
  h(t \mid \mathbf{x}_i) = h_0(t) \exp(\boldsymbol\beta^\top \mathbf{x}_i),
\]
where $h_0(t)$ is the unspecified baseline hazard and $\mathbf{x}_i$
is the covariate vector~\cite{cox1972}. Ties are handled via the
Breslow approximation. Model adequacy is assessed by the concordance
index~$C$, partial $\AIC$, and the proportional hazards assumption is
tested via scaled Schoenfeld residuals~\cite{grambsch1994}.

\paragraph{Accelerated failure time models.}
We fit parametric AFT models of the form
$\log T_i = \boldsymbol\mu^\top \mathbf{x}_i + \sigma\, W_i$,
where $W_i$ follows a standard extreme value (Weibull), normal
(log-normal), or logistic (log-logistic) distribution. Model
selection is by $\AIC$.

\paragraph{Competing risks.}
Resolution events are decomposed into proof ($k = 1$) and disproof
($k = 2$) as competing risks. We report cause-specific cumulative
incidence functions $\CIF_k(t) = P(T \le t, \text{cause} = k)$
and proved-to-disproved ratios stratified by tag.


%% ================================================================
\section{Results}
%% ================================================================

\subsection{Kaplan--Meier Estimates}

The overall Kaplan--Meier curve yields an estimated median survival
of $\tilde t = 946$ (i.e., the survival function crosses $0.50$
at problem number 946). The interquartile range is
$[\hat t_{0.25}, \hat t_{0.75}] = [570, 1134]$. At selected
timepoints: $\hat S(250) = 0.899$, $\hat S(500) = 0.782$,
$\hat S(750) = 0.654$, $\hat S(1000) = 0.451$.

Stratified Kaplan--Meier analysis reveals significant differences
by covariate group (Table~\ref{tab:logrank}).

\begin{table}[ht]
\centering
\caption{Log-rank tests for stratified Kaplan--Meier curves.}
\label{tab:logrank}
\begin{tabular}{@{}lcccl@{}}
\toprule
Stratum & $\tilde t_1$ & $\tilde t_0$ & $\chi^2$ & $p$-value \\
\midrule
Prize (yes/no)         & 703  & 965  & 24.29 & $< 0.001$ \\
Formalized (yes/no)    & 1102 & 900  & 24.22 & $< 0.001$ \\
OEIS (yes/no)          & ---  & ---  & 10.12 & $0.001$ \\
Graph theory (yes/no)  & 842  & ---  & 2.86  & $0.091$ \\
Ramsey theory (yes/no) & 894  & ---  & 0.01  & $0.906$ \\
\bottomrule
\end{tabular}
\raggedright
\small $\tilde t_1$, $\tilde t_0$: median survival for the indicated
and reference group, respectively. ``---'' indicates the median was
not reached or is omitted for brevity.
\end{table}

Prize problems have a median survival of 703, compared to 965 for
non-prize problems---a reduction of 262 ordinal units
($\chi^2_{LR} = 24.29$, $p < 0.001$). Conversely, formalized
problems exhibit a \emph{higher} median survival of 1102 versus 900,
which we discuss in Section~\ref{sec:discussion} as a selection effect.

\subsection{Cox Proportional Hazards Regression}

We fit a Cox model with structural covariates and subfield indicators
(Table~\ref{tab:cox}).

\begin{table}[ht]
\centering
\caption{Cox proportional hazards model with tag indicators
($n = 1{,}133$, events $= 483$).}
\label{tab:cox}
\begin{tabular}{@{}lcccc@{}}
\toprule
Covariate & $\hat\beta$ & $\HR$ & 95\% CI & $p$ \\
\midrule
Prize             & $0.96$  & $2.60$ & $(1.85, 3.67)$ & $< 0.001$ \\
Formalized        & $-0.59$ & $0.55$ & $(0.43, 0.71)$ & $< 0.001$ \\
OEIS              & $-0.47$ & $0.63$ & $(0.48, 0.82)$ & $< 0.001$ \\
Number theory     & $0.36$  & $1.44$ & $(1.13, 1.84)$ & $0.004$ \\
Graph theory      & $0.18$  & $1.20$ & $(0.93, 1.55)$ & $0.171$ \\
Ramsey theory     & $-0.18$ & $0.83$ & $(0.60, 1.16)$ & $0.284$ \\
Analysis          & $-0.12$ & $0.89$ & $(0.62, 1.26)$ & $0.503$ \\
\midrule
\multicolumn{5}{@{}l}{$C = 0.60$; partial $\AIC = 5838.0$;
LR $\chi^2 = 65.4$, $\mathrm{df} = 7$, $p < 10^{-10}$} \\
\bottomrule
\end{tabular}
\end{table}

Three covariates achieve significance at $\alpha = 0.01$: prize
attachment ($\HR = 2.60$; prize problems have $2.6\times$ the
instantaneous hazard of resolution), formalization
($\HR = 0.55$; 45\% reduction), and OEIS linkage
($\HR = 0.63$; 37\% reduction). Among tag indicators, only number
theory is significant ($\HR = 1.44$, $p = 0.004$), indicating that
number theory problems resolve faster than the baseline after
controlling for other covariates---a finding that may reflect the
larger community working in number theory rather than intrinsic
tractability.

Ramsey theory shows a non-significant tendency toward lower hazard
($\HR = 0.83$, $p = 0.284$), though the multivariate log-rank test
across five tag groups approaches significance
($\chi^2 = 8.63$, $\mathrm{df} = 4$, $p = 0.071$).

The global concordance index of $C = 0.60$ indicates modest but
real discriminative ability---consistent with the expectation that
problem-level metadata captures only a fraction of the variance in
resolution difficulty, the remainder being driven by the
(unobserved) intrinsic mathematical content.

\subsection{Competing Risks}

Decomposing resolution events into proof ($k = 1$) and disproof
($k = 2$), the overall proved-to-disproved ratio is $2.63{:}1$
($305$ proofs, $116$ disproofs), with 62 additional events classified
as ``solved'' without further specification. This ratio varies
substantially across subfields (Table~\ref{tab:competing}).

\begin{table}[ht]
\centering
\caption{Competing risks decomposition by mathematical subfield.}
\label{tab:competing}
\begin{tabular}{@{}lcccc@{}}
\toprule
Subfield & $n$ & Proved & Disproved & P\,:\,D ratio \\
\midrule
Overall         & 1133 & 305 & 116 & $2.6{:}1$ \\
Combinatorics   &   46 &  17 &   2 & $8.5{:}1$ \\
Ramsey theory   &  105 &  31 &   8 & $3.9{:}1$ \\
Analysis        &   74 &  29 &   9 & $3.2{:}1$ \\
Number theory   &  541 & 133 &  46 & $2.9{:}1$ \\
Geometry        &  101 &  25 &   9 & $2.8{:}1$ \\
Graph theory    &  253 &  74 &  35 & $2.1{:}1$ \\
Set theory      &   31 &   4 &   3 & $1.3{:}1$ \\
\bottomrule
\end{tabular}
\end{table}

The P\,:\,D ratio ranges from $8.5{:}1$ in combinatorics
to $1.3{:}1$ in set theory. Graph theory exhibits an elevated
disproof rate ($2.1{:}1$), consistent with the availability of
explicit counterexample constructions.

\subsection{Accelerated Failure Time Models}

We fit Weibull ($\AIC = 7793.7$), log-logistic ($7854.1$), and
log-normal ($7930.9$) AFT models; the Weibull provides the best fit
(Table~\ref{tab:aft}). The shape parameter $\hat\rho = 1.61$
(95\% CI: $1.49$--$1.74$) exceeds 1, indicating \emph{decreasing}
marginal hazard---problems surviving to higher indices have
progressively lower resolution rates.

\begin{table}[ht]
\centering
\caption{Weibull AFT model estimates ($n = 1{,}133$, events $= 483$).}
\label{tab:aft}
\begin{tabular}{@{}lcccc@{}}
\toprule
Covariate & $\hat\beta$ & $\AF = e^{\hat\beta}$ & 95\% CI & $p$ \\
\midrule
Prize      & $-0.51$ & $0.60$ & $(0.49, 0.74)$ & $< 0.001$ \\
Formalized & $0.32$  & $1.38$ & $(1.19, 1.61)$ & $< 0.001$ \\
OEIS       & $0.19$  & $1.21$ & $(1.03, 1.41)$ & $0.020$ \\
$n_{\text{tags}}$ & $-0.04$ & $0.96$ & $(0.87, 1.05)$ & $0.363$ \\
\midrule
\multicolumn{5}{@{}l}{Shape $\hat\rho = 1.61$ (95\% CI: $1.49$--$1.74$);
$\AIC = 7793.7$; $C = 0.58$} \\
\bottomrule
\end{tabular}
\raggedright
\small Acceleration factors $\AF < 1$ indicate faster resolution
(shorter ``survival''); $\AF > 1$ indicates slower resolution.
\end{table}

Prize attachment yields $\AF = 0.60$ (resolution in 60\% of the
time required by non-prize problems); formalization yields
$\AF = 1.38$ (38\% longer survival), reflecting the selection
effect discussed below.


%% ================================================================
\section{Discussion}
\label{sec:discussion}
%% ================================================================

\paragraph{The formalization paradox.}
Our most counterintuitive finding is that formalized problems
resolve \emph{more slowly} ($\HR = 0.55$, $\AF = 1.38$). This
is not a causal effect; rather, it is a classic \emph{selection bias}.
Problems chosen for formalization in Lean or Coq tend to be famous
and difficult---precisely those that resist resolution. The
formalization indicator acts as a proxy for unobserved difficulty:
only $21.7\%$ of formalized problems are resolved, versus $50.1\%$
of non-formalized problems.

\paragraph{The prize effect and Simpson's paradox.}
The prize hazard ratio of $\HR = 2.60$ admits two interpretations:
incentive effects attracting effort~\cite{kremer1998}, and more
precise problem statements facilitating targeted attacks.
Notably, the raw solve rate for prize problems ($38.0\%$) is
\emph{lower} than for non-prize problems ($41.2\%$)---the Cox model
reverses this after controlling for formalization and OEIS status,
which confound the association. This is Simpson's paradox in
survival data.

\paragraph{Competing risks and epistemic calibration.}
The $2.63{:}1$ proved-to-disproved ratio quantifies Erd\H{o}s's
\emph{epistemic calibration}: his conjectures were correct
$\sim\!70\%$ of the time among resolved problems. In combinatorics
($8.5{:}1$) his intuition was nearly infallible; in set theory
($1.3{:}1$) the near-parity suggests genuine uncertainty, consistent
with independence phenomena prevalent in that domain.

\paragraph{Limitations.}
(i)~Problem number is an ordinal proxy for time, not a metric variable;
the mapping from index to calendar time is nonlinear.
(ii)~The database lacks true resolution dates, preventing
calendar-time survival analysis.
(iii)~Covariates capture metadata rather than mathematical content;
the moderate concordance ($C = 0.58$--$0.60$) reflects this
inherent limitation.
Despite these constraints, the survival-analytic framework reveals
structure that descriptive statistics alone do not: Simpson's paradox
in the prize effect, the selection mechanism underlying the
formalization paradox, and field-dependent asymmetry in competing risks.

\paragraph{Implications.}
As AI-assisted theorem proving accelerates~\cite{trinh2024,
yang2024}, the resolution hazard may increase. Our Weibull baseline
provides a benchmark: if the post-2025 hazard exceeds the predicted
rate, this constitutes measurable evidence of AI impact on open
mathematical problems.


%% ================================================================
\paragraph{Acknowledgments.}
%% ================================================================
The Erd\H{o}s Problems database was compiled by Terence Tao and
contributors. Computational analyses used
\texttt{lifelines}~\cite{davidson-pilon2019} for survival modeling
and \texttt{scipy} for statistical tests. AI tools (Claude, Anthropic)
assisted with data analysis and manuscript preparation.


%% ================================================================
%% Bibliography
%% ================================================================
\begingroup
\raggedright
\begin{thebibliography}{99}

\bibitem{erdosproblems}
T.~Tao et al.,
``Erd\H{o}s Problems,''
\url{https://www.erdosproblems.com/}, 2024.

\bibitem{kaplan1958}
E.~L. Kaplan and P.~Meier,
``Nonparametric estimation from incomplete observations,''
\emph{J.\ Amer.\ Statist.\ Assoc.}, vol.~53, pp.~457--481, 1958.

\bibitem{cox1972}
D.~R. Cox,
``Regression models and life-tables (with discussion),''
\emph{J.\ Roy.\ Statist.\ Soc.\ Ser.~B}, vol.~34, pp.~187--220, 1972.

\bibitem{grambsch1994}
P.~M. Grambsch and T.~M. Therneau,
``Proportional hazards tests and diagnostics based on weighted residuals,''
\emph{Biometrika}, vol.~81, pp.~515--526, 1994.

\bibitem{davidson-pilon2019}
C.~Davidson-Pilon,
``lifelines: survival analysis in Python,''
\emph{J.\ Open Source Softw.}, vol.~4, no.~40, p.~1317, 2019.

\bibitem{borjas2015}
G.~J. Borjas and K.~B. Doran,
``Prizes and productivity: how winning the Fields Medal affects
scientific output,''
\emph{J.\ Human Resources}, vol.~50, pp.~728--758, 2015.

\bibitem{mihaljevi2019}
H.~Mihaljevi\'{c} and L.~Santamar\'{i}a,
``Mathematics publications and authors' gender:
learning from the Gender Gap in Science project,''
\emph{European J.\ Math.}, vol.~5, pp.~1--16, 2019.

\bibitem{cranshaw2011}
J.~Cranshaw and A.~Kittur,
``The Polymath Project: lessons from a successful online
collaboration in mathematics,''
in \emph{Proc.\ CHI}, ACM, 2011, pp.~1865--1874.

\bibitem{kremer1998}
M.~Kremer,
``Patent buyouts: a mechanism for encouraging innovation,''
\emph{Quarterly J.\ Economics}, vol.~113, pp.~1137--1167, 1998.

\bibitem{trinh2024}
T.~H. Trinh et al.,
``Solving olympiad geometry without human demonstrations,''
\emph{Nature}, vol.~625, pp.~476--482, 2024.

\bibitem{yang2024}
K.~Yang et al.,
``LeanDojo: theorem proving with retrieval-augmented language models,''
in \emph{NeurIPS}, 2024.

\end{thebibliography}
\endgroup

\end{document}
